\documentclass[12pt]{article}

\usepackage[T1]{fontenc}
\usepackage[utf8]{inputenc}
\usepackage{lmodern}
\usepackage[margin=1in]{geometry}
\usepackage{amsmath,amssymb}
\usepackage{siunitx}
\usepackage{microtype}
\usepackage{booktabs}
\usepackage{enumitem}
\usepackage[hidelinks]{hyperref}

\sisetup{
  detect-all,
  separate-uncertainty=true,
  range-phrase=--,
  range-units=single
}

\newcommand{\placeholdercite}[1]{\textbf{[Citation needed: #1]}}

\title{Falcon 9 Seismo-acoustic Analysis: Methods Draft}
\author{}
\date{Draft reconstructed from the canonical analysis notebooks}

\begin{document}
\maketitle

\section{Methods}

\subsection{Analysis overview}

The analysis was designed to reconstruct the temporal evolution of the Falcon 9
accident using synchronized infrasound and seismic observations from the BCHH
array. The workflow comprised: (1) correction of the instrumental responses and
conversion to physical units; (2) compilation of surveyed source--receiver
geometry and meteorological constraints; (3) manual identification and
association of candidate airwave arrivals; (4) measurement of acoustic pulse
morphology and signal-to-noise ratio; (5) estimation of back azimuth and apparent
propagation speed using three-channel planar-wave coherence analysis; (6)
comparison with a finite-range circular-wavefront model; and (7) measurement of
acoustic pressure, seismic peak ground velocity, and their amplitude
relationship.

All quantitative processing was performed on continuously recorded waveforms
sampled at \SI{250}{\hertz}. The canonical analysis interval extended from
\SI{180}{\second} before to \SI{1800}{\second} after the initial upper-stage
explosion at 13:07:12.080 UTC.

\subsection{Instrument calibration and response correction}

The three seismic components were corrected using the manufacturer response for
the Nanometrics Trillium Compact Posthole seismometer--Centaur digitizer
combination. The combined sensitivity was
$3.0172\times10^{8}$ counts~$(\si{\metre\per\second})^{-1}$. The Centaur
digitizer was operated on its \SI{40}{\volt} peak-to-peak input range,
corresponding to $4.0\times10^{5}$ counts~\si{\per\volt} and an implied
seismometer sensitivity of approximately
\SI{754}{\volt\per\metre\second}.

The three infraBSU differential-pressure sensors were converted independently
from counts to pressure using empirical, sensor-specific field calibrations. The
adopted sensitivities were 8.10, 0.690, and 10.0 counts~\si{\per\pascal} for HD1,
HD2, and HD3, respectively. These factors were derived from comparative
calibration experiments in which co-located infraBSU sensors were exposed to
common pressure changes and impulsive signals. The calibrations should therefore
be regarded as empirical field estimates rather than laboratory-traceable
absolute calibrations.

For the infrasound channels, the median count level between 20 and \SI{2}{\second}
before the initial explosion was removed before response correction. The
empirical sensitivities were incorporated into complete instrument-response
representations while retaining the nominal poles, zeros, and digitizer stages.
Seismic traces were demeaned and linearly detrended. All traces received a 5\%
cosine taper before frequency-domain response removal.

A four-corner prefilter of 0.02, 0.05, 80, and \SI{100}{\hertz} and a water level
of \SI{60}{\decibel} were used to stabilize the response inversion. This
prefilter was applied only as part of the response-removal operation; no
additional Butterworth filtering was applied during event association, amplitude
measurement, or array analysis. The corrected seismic channels are expressed as
ground velocity in \si{\metre\per\second}, and the corrected infrasound channels
as differential pressure in \si{\pascal}.

\subsection{Source--receiver geometry}

Sensor coordinates were obtained from the differential-GPS survey of BCHH and
transformed to UTM Zone 17N for array calculations. Horizontal source--receiver
distances and geographic azimuths were independently calculated on the WGS84
ellipsoid. The array reference position was the central broadband seismometer at
approximately \SI{541820.4}{\metre} E and \SI{3160866.5}{\metre} N. The reference
distance between SLC-40 and BCHH was \SI{1417.9}{\metre}, and the back azimuth
from the array toward SLC-40 was \ang{199.4}, measured clockwise from north.

HD2 was used as the infrasound reference channel. The individual
source-to-sensor distances were retained when calculating differential acoustic
travel times. This avoided treating the approximately \SI{30}{\metre} array as a
single receiver during pick association and waveform alignment.

Approximate source times assigned to the major accident stages were obtained
from the video chronology. These times were used to label and interpret key
signals but were not required for constructing the general candidate-event
catalogue, which was based on the observed infrasound arrivals.

\subsection{Meteorological correction}

Meteorological observations from the Kennedy Space Center tower network between
13:05 and 13:35 UTC were used to estimate the effective acoustic propagation
speed. The selected dataset contained 84 temperature and 84 relative-humidity
observations. Mean conditions were approximately \SI{26.36}{\degreeCelsius} and
83.87\% relative humidity.

The still-air sound speed was calculated as
\begin{equation}
  c_0 = 331.3 + 0.606T + 1.26h,
  \label{eq:still-air-sound-speed}
\end{equation}
where $T$ is temperature in degrees Celsius and $h$ is relative humidity
expressed as a fraction. This yielded
$c_0=\SI{348.33}{\metre\per\second}$.

Wind observations were resolved into east and north components, averaged within
\SI{10}{\metre} height intervals, and linearly interpolated over the lowest
\SI{65}{\metre} of the atmosphere. The height-averaged wind component parallel
to the SLC-40--BCHH propagation path was then added to the still-air sound speed:
\begin{equation}
  c_{\mathrm{eff}} = c_0 + \mathbf{u}\cdot\hat{\mathbf{r}}.
  \label{eq:effective-sound-speed}
\end{equation}
The mean along-path wind component was
$+\SI{2.65}{\metre\per\second}$, giving an effective acoustic speed of
\SI{350.98}{\metre\per\second}. The corresponding travel time over
\SI{1417.9}{\metre} was \SI{4.040}{\second}, approximately
\SI{0.031}{\second} shorter than the still-air estimate.

\subsection{Baseline removal and manual arrival picks}

Candidate airwave arrivals were initially identified by visual inspection using
Antelope \texttt{dbpick}. Accepted picks were restricted to non-deleted phase-N
arrivals on HD1, HD2, and HD3. The original pick precision was retained during
export to the Python analysis workflow.

Slow variations in the infrasound baseline were removed with a centered,
noncausal moving median. For each infrasound trace, a 251-sample window---
approximately \SI{1}{\second} at \SI{250}{\hertz}---was used to estimate the
baseline, which was then subtracted from the response-corrected pressure
waveform. The three seismic channels were retained without this baseline
operation. Because a centered median is nonlinear, it is not formally a
zero-phase filter, but it introduces no systematic delay equivalent to that of a
one-sided causal smoother.

The moving-median window was chosen after comparing unfiltered waveforms,
high-pass-filtered waveforms, local linear baselines, and moving-median windows
of several durations. High-pass filtering improved some weak arrivals but
degraded broader pulses and altered their morphology. Local linear baselines
were sometimes controlled by a small number of samples and introduced
implausibly steep trends. The \SI{1}{\second} moving median provided the most
stable baseline while preserving the short pressure impulses.

\subsection{Reduced-time association of candidate events}

Before association, each accepted pick was corrected for its expected
differential travel time relative to HD2:
\begin{equation}
  t_{i,\mathrm{red}}
  = t_i - \frac{r_i-r_{\mathrm{HD2}}}{c_{\mathrm{eff}}},
  \label{eq:reduced-pick-time}
\end{equation}
where $r_i$ is the SLC-40--sensor distance and
$c_{\mathrm{eff}}=\SI{350.98}{\metre\per\second}$. The resulting corrections
were approximately \SI{0.090}{\second} for HD1, \SI{0}{\second} for HD2, and
\SI{0.013}{\second} for HD3. These corrections were used only to associate
likely corresponding arrivals; the final propagation direction and speed were
estimated independently.

Reduced picks were sorted by time. Beginning with the earliest unassigned pick,
all picks within the following \SI{0.040}{\second} were considered as one
candidate cluster. At most one pick per channel was retained, selecting the pick
closest to the cluster seed when multiple picks occurred on a channel. A cluster
was accepted as a candidate event when it contained picks from at least two
distinct infrasound channels. Picks that did not satisfy this criterion remained
unassociated rather than being forced into an event.

A fixed six-channel waveform segment extending \SI{0.20}{\second} before to
\SI{0.30}{\second} after each reduced event time was saved for subsequent
analysis. The provisional time shifts were not written into these event streams;
their original observed sample times were preserved.

\subsection{Acoustic pulse measurements and candidate tiers}

For initial pulse characterization, the three infrasound traces were shifted to
the HD2 reference time using the meteorologically predicted differential delays
and combined using a sample-by-sample median stack. A positive compression
followed by a negative excursion was sought between $-0.04$ and
$+\SI{0.25}{\second}$ relative to the reduced event time. The negative peak was
required to occur between 0.010 and \SI{0.180}{\second} after the positive peak.
When multiple pairs satisfied these conditions, the pair with the largest
peak-to-peak pressure was selected.

For a positive peak at $t_+$ and a negative peak at $t_-$, provisional one-cycle
bounds were defined as
\begin{align}
  t_{\mathrm{start}} &= t_+ - 0.5(t_- - t_+), \\
  t_{\mathrm{end}}   &= t_- + 0.5(t_- - t_+).
  \label{eq:event-bounds}
\end{align}

Noise was characterized over $-0.40$ to $-\SI{0.10}{\second}$ using the robust
scale
\begin{equation}
  \sigma_{\mathrm{robust}}
  = 1.4826\,\operatorname{median}
  \left(\left|p-\operatorname{median}(p)\right|\right).
  \label{eq:robust-noise}
\end{equation}
The initial stack SNR was the maximum absolute pressure within the event bounds
divided by this noise scale. The positive--negative morphology was used as a
measurement-quality diagnostic, not as the fundamental definition of an event.
The principal explosion, for example, was retained despite failing the simple
one-cycle morphology because it was an exceptionally strong, coherent
multi-channel signal with a more complex pressure history.

Candidate labels based on channel support, reduced-pick span, and SNR were
retained as descriptive quality indicators. They were not treated as distinct
physical classes.

\subsection{Eligibility for array analysis}

Two configurable threshold sets were applied. The permissive analysis set
required a median-stack peak-to-peak pressure of at least \SI{10}{\pascal}, stack
SNR of at least 1.5, support from at least two channels, and a valid
positive--negative measurement. The named principal explosion was explicitly
retained despite failing the short-pulse morphology criterion.

A more conservative core subset required at least \SI{25}{\pascal} peak-to-peak
pressure and SNR of at least 3, with the same channel-support requirement. These
thresholds were used to distinguish the most readily interpretable events from
additional candidates; they do not represent sharp natural boundaries between
signal and noise.

\subsection{Planar-wave array analysis}

Back azimuth and apparent horizontal propagation speed were estimated jointly
by shifting all three infrasound waveforms according to trial plane-wave models.
Back azimuth was defined as the direction from the array toward the source,
clockwise from north. For a trial back azimuth $\theta$ and speed $v$, the
propagation vector toward the array was
\begin{equation}
  \mathbf{k} = [-\sin\theta,\,-\cos\theta],
  \label{eq:propagation-vector}
\end{equation}
and the predicted delay at sensor $i$, relative to HD2, was
\begin{equation}
  \Delta t_i =
  \frac{(\mathbf{x}_i-\mathbf{x}_{\mathrm{HD2}})\cdot\mathbf{k}}{v}.
  \label{eq:planar-delay}
\end{equation}

For each trial, the three traces were interpolated after applying the predicted
delays and trimmed to the same common overlap. Thus, every point on the search
surface was evaluated using the same waveform duration, avoiding the
unequal-overlap bias that can favor zero lag in conventional cross-correlation.

Each aligned trace was median-centered and normalized by its RMS amplitude.
Coherence was quantified using both three-channel semblance,
\begin{equation}
  S =
  \frac{\sum_t\left[\sum_i x_i(t)\right]^2}
  {N\sum_t\sum_i x_i^2(t)},
  \label{eq:semblance}
\end{equation}
and the mean of the three pairwise correlation coefficients. The pairwise mean
was mapped from $[-1,1]$ to $[0,1]$, and the joint score was
\begin{equation}
  C = 0.60S + 0.40\left(\frac{\overline{\rho}+1}{2}\right).
  \label{eq:joint-coherence}
\end{equation}

The coarse search covered back azimuths from \ang{180} to \ang{220} in
\ang{0.5} increments and speeds from 280 to \SI{440}{\metre\per\second} in
\SI{2}{\metre\per\second} increments. Each coarse maximum was refined over
$\pm\ang{2}$ in \ang{0.1} increments and
$\pm\SI{20}{\metre\per\second}$ in \SI{0.5}{\metre\per\second} increments.
Most events were scored over $-0.04$ to $+\SI{0.18}{\second}$; the principal
explosion used an extended interval ending at $+\SI{0.35}{\second}$.

Pairwise correlations and residual lags were calculated after applying the
best-fitting model. Lag closure,
\begin{equation}
  \tau_{12}+\tau_{23}+\tau_{31},
  \label{eq:lag-closure}
\end{equation}
was retained as a consistency diagnostic. The relative performance of HD1 was
described by comparing the HD2--HD3 correlation with the mean of the two pairs
containing HD1.

Solution stability was tested by shifting the scoring window by
$-0.008$, 0, and $+\SI{0.008}{\second}$ and scaling its duration by 0.9, 1.0,
and 1.1. The standard deviations of the resulting back azimuths and speeds were
retained as empirical stability measures. Near-maximum widths at score
reductions of 0.01 and 0.02 were also calculated, but these widths are
descriptive properties of the coherence surface and not formal confidence
intervals.

Composite coherence surfaces were constructed by combining the coarse surfaces
for all analyzed events and for the core subset. Arithmetic-mean, median,
maximum-score-weighted, and individually normalized mean composites were
retained.

\subsection{Finite-range circular-wavefront analysis}

The plane-wave approximation was evaluated using exact source--receiver
distances in the horizontal plane. For a trial source coordinate $\mathbf{x}_s$
and propagation speed $v$, the predicted relative delay was
\begin{equation}
  \Delta t_i =
  \frac{
    \left|\mathbf{x}_i-\mathbf{x}_s\right|
    -\left|\mathbf{x}_{\mathrm{HD2}}-\mathbf{x}_s\right|
  }{v}.
  \label{eq:circular-delay}
\end{equation}
Because receiver elevations were not available in the adopted geometry, this is
strictly a two-dimensional circular-wavefront model rather than a full
three-dimensional spherical-wave model.

The primary finite-range calculation held the source at SLC-40 and searched
speeds from 280 to \SI{440}{\metre\per\second} in
\SI{0.5}{\metre\per\second} increments. Coherence was calculated with the same
waveform windows and scoring metric used for the planar solutions.

Source-position sensitivity was examined in coordinates oriented along and
perpendicular to the SLC-40--BCHH path. Along-range offsets from $-500$ to
$+\SI{500}{\metre}$ and cross-range offsets from $-150$ to
$+\SI{150}{\metre}$ were evaluated while profiling over the complementary
coordinate. These calculations were treated as resolution tests rather than
source locations. Three sensors provide only two independent relative delays, so
source easting, source northing, and propagation speed cannot be determined
independently from these data. Open along-range coherence ridges were therefore
interpreted as a fundamental range-resolution limitation.

\subsection{Acoustic and seismic amplitude measurements}

Acoustic and seismic amplitudes were measured for every candidate event using
the saved baseline-corrected waveform segments. Event intervals were taken from
the positive--negative pulse measurements, except for the principal explosion,
for which a fixed $-0.04$ to $+\SI{0.35}{\second}$ interval was used.

For each seismic component, the median of the pre-event interval from $-0.18$ to
$-\SI{0.05}{\second}$ was removed. Component PGV was the largest absolute ground
velocity within the event interval. Three-component vector PGV was calculated as
\begin{equation}
  \mathrm{PGV}_{\mathrm{vec}}
  = \max_t\sqrt{v_E^2(t)+v_N^2(t)+v_Z^2(t)}.
  \label{eq:vector-pgv}
\end{equation}
The maximum vector amplitude over the entire saved segment was also retained to
test whether the acoustic event interval captured the largest seismic motion.

The infrasound traces were aligned to HD2 using the meteorological differential
delays and combined as a median stack. Positive, negative, and peak-to-peak
pressures were measured within the same event interval. Acoustic SNR was defined
as stack peak-to-peak pressure divided by the robust pre-event noise scale.

Where pressure was reduced to a reference distance of \SI{1}{\kilo\metre},
spherical geometric spreading was assumed:
\begin{equation}
  p_{\mathrm{red},\,1\,\mathrm{km}}
  = p_{\mathrm{obs}}\frac{r}{\SI{1000}{\metre}}.
  \label{eq:reduced-pressure}
\end{equation}
Peak sound-pressure level was calculated from the positive pressure peak as
\begin{equation}
  L_{\mathrm{peak}}
  = 20\log_{10}\left(\frac{p_+}{\SI{20}{\micro\pascal}}\right).
  \label{eq:peak-spl}
\end{equation}
This quantity is explicitly a peak pressure level rather than an RMS
sound-pressure level.

The dimensional acoustic-to-seismic amplitude ratio was calculated as
peak-to-peak pressure divided by vector PGV, with units of
\si{\pascal\per\metre\second}. It was used as an empirical amplitude-coupling
measure and should not be interpreted as an acoustic-to-seismic energy ratio.

\subsection{Acoustic--seismic amplitude regression}

Regression was restricted to events with seismic vector-PGV SNR of at least 5,
acoustic peak-to-peak SNR of at least 5, and an event-window PGV at least 70\% of
the maximum PGV in the complete saved segment.

A free power-law model,
\begin{equation}
  p_{\mathrm{p2p}} = A\,\mathrm{PGV}^{b},
  \label{eq:free-power-law}
\end{equation}
was fitted by ordinary least squares after taking base-10 logarithms. A 95\%
descriptive interval for the exponent was estimated using 5000 nonparametric
event bootstrap resamples. A proportional model,
\begin{equation}
  p_{\mathrm{p2p}} = K\,\mathrm{PGV},
  \label{eq:proportional-model}
\end{equation}
was also evaluated as a constant-ratio alternative. The two models were compared
using root-mean-square residuals in base-10 logarithmic space. These fits
describe the observed relationship and do not account for uncertainty in both
axes, calibration covariance, or censoring near the noise floor.

\subsection{Video and audio comparison}

Public video recordings were used to identify major visible stages of the
accident and to assign approximate source-time labels to the upper-stage
explosion, principal lower-stage explosion, capsule impact, and subsequent
capsule-related explosion. The current quantitative catalogue and array results
remain derived from the BCHH geophysical data.

For the supplementary opening-sequence visualization, audio extracted from one
YouTube recording was demeaned and band-pass filtered from 1 to \SI{2}{\kilo\hertz}
using a four-pole zero-phase filter. Because the video did not contain an
absolute timestamp, its timing was aligned empirically with the BCHH sequence.
The audio timing and source-reduced overview should therefore be treated as
illustrative synchronization rather than an independent absolute timing
observation.

\subsection{Reproducibility}

Each processing stage writes its numerical products, configuration parameters,
intermediate event streams, and diagnostic figures to a dedicated output
directory. The canonical workflow comprises the response-correction,
input-preparation, meteorological, candidate-association, planar-array,
circular-wavefront, acoustic--seismic amplitude, and paper-product notebooks.
Event-filter decisions are retained for every candidate, including events that
fail the permissive or core thresholds.

\clearpage
\section*{Items requiring confirmation before submission}

\begin{enumerate}[leftmargin=*]
  \item \textbf{Calibration provenance and uncertainty.} Verify the exact mix
  of lift tests, door-slam experiments, and other impulsive tests used to derive
  the BCHH factors. The proposed $\pm20\%$ pressure uncertainty needs a documented
  basis.

  \item \textbf{Principal-explosion pressure definition.} The current
  median-stack measurement and the earlier sensor-specific extrema give
  different positive and negative pressures. Select one consistent headline
  measure and identify the other explicitly as a different estimator.

  \item \textbf{Back azimuth to SLC-40.} Freeze one geometry file and use its
  value consistently; current and archived products contain values between
  approximately \ang{199.0} and \ang{199.9}.

  \item \textbf{Acoustic speeds.} Do not conflate the meteorologically derived
  \SI{350.98}{\metre\per\second} with the approximately
  \SIrange{366.5}{366.8}{\metre\per\second} empirical display alignment used in
  the video/audio overview.

  \item \textbf{Event boundaries.} The canonical measurement presently uses
  the half-lag extension around the positive and negative peaks. A configuration
  flag for return-to-noise refinement is present but unused.

  \item \textbf{Event-filter thresholds.} Describe the permissive and core
  thresholds as operational divisions along continuous distributions, not as
  natural signal/noise discontinuities.

  \item \textbf{Wavefront terminology.} With horizontal geometry only, the
  adopted finite-range calculation is circular rather than fully spherical.

  \item \textbf{Video analysis.} Add camera locations, video identifiers, frame
  rates, synchronization procedures, and timing uncertainties before claiming a
  quantitative frame-by-frame analysis.

  \item \textbf{Energy estimates.} Acoustic and seismic energies and equivalent
  magnitudes are not regenerated by the current canonical workflow. Omit them
  until a reproducible analysis product is added.

  \item \textbf{Source-time uncertainties.} Assign and justify uncertainties
  for the four named source-time markers.
\end{enumerate}

\section*{Citation checklist}

Add references for Antelope and \texttt{dbpick}; ObsPy response removal; the
relevant NumPy, SciPy, pandas, Matplotlib, and pyproj software; the sound-speed
temperature--humidity approximation; the KSC meteorological data; semblance and
array-slowness analysis; the Trillium Compact Posthole, Centaur, and infraBSU
responses; robust MAD scaling; nonparametric bootstrap methods; and each public
video used in the chronology.

\end{document}
